\documentclass[10pt,a4paper]{amsart}
\usepackage[margin=19mm]{geometry}
\usepackage{amsmath,amssymb,mathtools}
\usepackage[scr=rsfs,cal=euler]{mathalfa}
\usepackage{xfrac}
\usepackage[unicode,hidelinks,bookmarks=false]{hyperref}
\newcommand{\ie}{i.e.\ }
\newcommand{\cf}{cf.\ }
\newcommand{\resp}{respectively\ }
\newcommand{\Subsection}{\subsection}
\providecommand{\nobreakdash}{\nobreak\hbox{-}}
\setlength{\emergencystretch}{2em}
\multlinegap0pt
\usepackage{bm}
\usepackage{bbm}
\usepackage{dsfont}
\usepackage{xspace}
\usepackage{cancel}
\usepackage{mathtools}
\usepackage{amsthm}
\makeatletter
\@ifundefined{dfn}{\newtheorem{dfn}{Dfn}}{}
\@ifundefined{lem}{\newtheorem{lem}[dfn]{Lemma}}{}
\makeatother
\newcommand{\R}{\mathbb R}
\newcommand{\ma}{\measuredangle}
\title[A Splitting Theorem for non-positively curved Lorentzian spa]{A Splitting Theorem for non-positively curved Lorentzian spaces}
\author{Joe Barton, Tobias Beran, Mauricio Che, Sebastian Gieger, Jona R\"ohrig, Felix Rott}
\date{Source: arXiv:2601.14058v1 (CC BY 4.0)}
\begin{document}
\maketitle

\noindent\emph{Source and licence.} A Splitting Theorem for non-positively curved Lorentzian spaces. Version arXiv:2601.14058v1 (CC BY 4.0), distributed under the Creative Commons Attribution 4.0 International licence, as recorded in the arXiv licence metadata for this exact version and shown on the abstract page https://arxiv.org/abs/2601.14058v1. Full rights notice: licenses/arxiv-2601.14058-CC-BY-4.0.txt.\par\smallskip

\noindent\textit{Editorial scope.} Counted unit: the Lem whose source label is \texttt{lem:firstvargeq} in the article (source0.tex lines 506--517, source label \texttt{lem:firstvargeq}), delivered together with the article's own complete original proof of it (source0.tex lines 519--562, one \texttt{proof} environment of 44 lines). No mathematics is written, repaired or re-derived: statement and proof are the article's own text line for line. Required context retained uncounted: none. Every definition, display and label of those lines is retained unchanged. Omitted from this package and not redistributed: 1--505, 518--518, 563--1953. Nothing inside a retained excerpt is omitted or altered: every retained body is delivered exactly as the article's lines are written, so no omission receipt of any kind accompanies this package. Citations: the retained excerpts carry no citation command at all, so no bibliography entry of the article is transcribed and no reference is redistributed. Reference adaptation: none was needed and none was made; every reference inside the delivered unit resolves inside it, so no unresolved reference or citation is produced. Printed slips kept literal: no printed word, symbol, hypothesis or number of the counted statement or of its proof was altered. Layout and class: the article's own class is \texttt{amsart} with packages including inputenc, amssymb, amsmath, amsthm, amsfonts, amscd, mathtools, ulem, hyperref, zref-clever, enumitem, xcolor, microtype, tikz; the wrapper is the lane's pinned \texttt{amsart} editorial wrapper at 10pt on A4, which restates the theorem-like environments the retained text needs and adds the article's own macro definitions that the retained text needs (\texttt{\textbackslash R}, \texttt{\textbackslash ma}), with extra wrapper packages (bm, bbm, dsfont, xspace, cancel, mathtools). The delivered numbering is the unit's own. Assets: no retained span contains a figure environment, a graphics-inclusion command, a PDF-page inclusion, a table or a TikZ picture; no image file is copied into this package, the package declares no asset dependency and no image file is redistributed with it. Licence: version arXiv:2601.14058v1 of this article is distributed under the Creative Commons Attribution 4.0 International licence, as recorded in the arXiv licence metadata for this exact version and shown on the abstract page https://arxiv.org/abs/2601.14058v1. A full rights notice is delivered at licenses/arxiv-2601.14058-CC-BY-4.0.txt. Characters: every retained excerpt is pure ASCII, so the delivered engine, XeLaTeX with its default Latin Modern text font, reports no missing character.
\begin{lem}[First variation formula, second inequality in the model space]
\label{lem:firstvargeq}
     Let $K\in\mathbb R$ and $\triangle(a, b, c)$ be a timelike triangle in $\mathbb L^2(K)$ satisfying size bounds such that either $b\ll a\ll c$ or $a\ll c\ll b$. Moreover, let $\theta:=\ma^S_a(b,c)$. Calling the side lengths of the triangle $t:=\tau_s(a,c),\ z:=\tau_s(b,c)$ and $y:=\tau_s(b,a)$, we obtain
    \begin{equation*}
        \sigma\cosh(\theta)\geq\frac{z-y}{t} \, ,
    \end{equation*}
    where 
    \[\sigma:=\begin{cases}
    +1 & b\ll a\ll c,\\
    -1 & a\ll c\ll b.
\end{cases}\ \]
\end{lem}

\begin{proof}
    We first show the proof for $K=0$ and then for general $K\in\R$.
    By the Lorentzian law of cosines, we have
    \begin{align*}
        \sigma \cosh(\theta)&=\frac{z^2-y^2-t^2}{2yt}\\
        &=\frac{z-y}{t} \frac{z+y}{2y}-\frac{t}{2y} - \frac{z-y}{t} + \frac{z-y}{t}\\
        &=\underbrace{\sigma\frac{z-y}{t}}_{\geq1}\underbrace{\sigma\left(\frac{z+y}{2y}-1\right)}_{=\sigma\frac{z-y}{2y}\geq\frac{t}{2y}}-\frac{t}{2y}+\frac{z-y}{t}\\
        &\geq\frac{t}{2y}-\frac{t}{2y}+\frac{z-y}{t}=\frac{z-y}{t} \, ,
    \end{align*}
    where we used multiple times that $\sigma(z-y)\geq t$ by the reverse triangle inequality. \\

    For $K=1$ and $\sigma=+1$, we must prove
    \begin{equation*}
        \frac{z-y}{t} \leq 
        \sigma \cosh(\theta) =\frac{\cosh(z) - \cosh(t)\cosh(y)}{\sinh(t) \sinh(y)} 
        =\frac{\cosh(z) - \cosh(t+y)}{\sinh(t) \sinh(y)} +1 \, , 
    \end{equation*}
    where we have used the addition theorem for $\cosh(t+y)$. By the reverse triangle inequality, $z-(t+y)\geq 0$, and hence, the above is equivalent to 
    \begin{equation}\label{eq:equivalentformulation}
    \frac{\cosh(z)-\cosh(t+y)}{z-(t+y)} \geq \frac{\sinh(t)\sinh(y)}{t} \, .
    \end{equation}
    We observe that the left hand side is a difference quotient of $\cosh$. The mean value theorem states that there is a $\lambda\in[t+y,z]$ such that 
    \[\sinh(\lambda)=\frac{d}{d s}\bigg|_{s=\lambda}\cosh(s) = \frac{\cosh(z)-\cosh(t+y)}{z-(t+y)} \, .\]
    Since $\sinh$ is monotonically increasing in $\R$ and $\lambda\geq t+y$, to prove \eqref{eq:equivalentformulation}, it is enough to show
    \begin{equation}\label{eq:almostequivalent}
    \sinh(t+y) \geq \frac{\sinh(t)\sinh(y)}{t} \, .
    \end{equation}
    By the addition theorem for $\sinh(t+y)$ and the fact that for $t>0$ $\frac{\sinh(t)}{t}\leq \cosh(t)$, we see
    \[\sinh(t+y)=\sinh(t)\cosh(y)+\cosh(t)\sinh(y)\geq\frac{\sinh(t)\sinh(y)}{t}\]
    proving \eqref{eq:almostequivalent} and subsequently \eqref{eq:equivalentformulation}. 

    In fact, the equality holds for arbitrary $K>0$ by substituting $\sqrt{K}t$, $\sqrt{K}y$ and $\sqrt{K}z$ for $t$, $y$ and $z$ respectively:
    \[ \sigma \cosh(\theta) =\frac{\cosh(\sqrt{K}z) - \cosh(\sqrt{K} t)\cosh(\sqrt{K} y)}{\sinh(\sqrt{K} t) \sinh(\sqrt{K} y)} \geq \frac{\sqrt{K} z - \sqrt{K} y}{\sqrt{K} t} = \frac{z-y}{t} \, .\]

    The other three cases ($K=\pm1, \sigma=\pm1$) work similarly by using different inequalities.
    Instead of repeating the computations, we only mention the differences to the previous case. 
    
    For $K=1,\sigma=-1$, we initially use the addition theorem for $\cosh(y-t)$ instead of $\cosh(y+t)$ and apply the reverse triangle inequality $y\geq t+z$. At the end, instead of $\frac{\sinh(t)}{t}\leq\cosh(t)$, we use $\frac{\sinh(t)}{t}\geq 1$, which then yields the same result. 
    In the case $K=-1$, the inequality we need to prove reads 
    \[\frac{\cos(t)\cos(y) - \cos(z)}{\sin(t) \sin(y)} \geq \frac{z-y}{t} \, .\]
    We again distinguish the cases $\sigma=\pm1$. At the beginning, if $\sigma=+1$, we apply the addition theorem to $\cos(t+y)$ while for $\sigma=-1$ we apply it to $\cos(y-t)$. At the end, if $\sigma=+1$, we use that $\frac{\sin(t)}{t}\leq 1$ and if $\sigma=-1$, we use that $\frac{\sin(t)}{t}\geq \cos(t)$ for $t\in[0,\pi]$.  

    Finally, for arbitrary $K\in\R$, we note that in the case $K>0$ and $\sigma=+1$, the inequalities do not change when scaling all the side lengths by $\sqrt{|K|}$, which finishes the proof.
\end{proof}

\end{document}
